\subsection{The mass bounds as entropy differences}

In a closed background the difference between $H_0$ and the entropy is constant. In general it measures how much the total Monge--Amp\`ere mass of bounded potentials varies: its infimum and supremum recover the lower and upper masses $v_-(\omega)$ and $v_+(\omega)$ of \cref{def:bmp}.

Let $\mathcal P_{\mathrm{ent}}(X)$ denote the set of probability measures of finite entropy relative to the fixed smooth probability volume $\mathrm{d}V$. By \cref{prop:hermitian-dual-entropy-bounds}, the difference $H_0(\nu)-\operatorname{Ent}_{\mathrm{d}V}(\nu)$ is a finite real number for every $\nu$ in this set.

\Needspace{13\baselineskip}
\begin{proposition}\label{prop:entropy-mass-extrema}\index{Monge--Ampere mass@Monge--Amp\`ere mass!entropy characterization}
For every Hermitian background $\omega$,
\[
 \begin{split}
 \inf_{\nu\in\mathcal P_{\mathrm{ent}}(X)}
      \bigl(H_0(\nu)-\operatorname{Ent}_{\mathrm{d}V}(\nu)\bigr)
       &=\log v_-(\omega),\\
 \sup_{\nu\in\mathcal P_{\mathrm{ent}}(X)}
      \bigl(H_0(\nu)-\operatorname{Ent}_{\mathrm{d}V}(\nu)\bigr)
       &=\log v_+(\omega).
 \end{split}
\]
The same extrema are obtained over smooth strictly positive probability measures. Here $\log0=-\infty$ and $\log(+\infty)=+\infty$. Thus the difference is bounded above if and only if $X$ has the bounded mass property, and bounded below if and only if $\omega$ has positive lower mass.
\end{proposition}
\begin{proof}
Write $Z_G=\int_X\mathrm{e}^G\,\mathrm{d}V$. A bounded solution for $G\in C(X)$ has mass $\mathrm{e}^{B_0(G)}Z_G$. If $v_-(\omega)>0$, the lower mass bound gives
\[
 B_0(G)\geq\log v_-(\omega)-\log Z_G.
\]
Taking the defining supremum for $H_0$ and using \cref{lem:entropy-variational} gives $H_0-\operatorname{Ent}_{\mathrm{d}V}\geq\log v_-(\omega)$ on $\mathcal P_{\mathrm{ent}}(X)$. In the same way, if $v_+(\omega)<\infty$, the upper mass bound gives $H_0-\operatorname{Ent}_{\mathrm{d}V}\leq\log v_+(\omega)$. These bounds hold trivially when $v_-(\omega)=0$ or $v_+(\omega)=+\infty$.

To prove sharpness, first take a smooth real $G$ and its Gibbs probability $\nu_G=Z_G^{-1}\mathrm{e}^G\,\mathrm{d}V$, which has finite entropy. Testing the supremum defining $H_0$ with $G$ gives
\[
 H_0(\nu_G)-\operatorname{Ent}_{\mathrm{d}V}(\nu_G)
       \geq B_0(G)+\log Z_G.
\]
In the other direction, let $\sigma_G$ be the smooth strictly positive Gauduchon adjoint probability of \cref{thm:normalization-gauduchon-recovery}. It lies in $\mathcal A_0(G)$, so its dual value is $B_0(G)+\int_XG\,\mathrm{d}\sigma_G$, and \cref{lem:entropy-variational} yields
\[
 H_0(\sigma_G)-\operatorname{Ent}_{\mathrm{d}V}(\sigma_G)
       \leq B_0(G)+\log Z_G.
\]
As $G$ runs over smooth real functions, the masses $\mathrm{e}^{B_0(G)}Z_G$ run over the masses $\int_X(\omega+\ddc u)^n$ of all smooth $u$ with $\omega+\ddc u>0$, since such a $u$ solves the equation for $G=\log((\omega+\ddc u)^n/\mathrm{d}V)$. By \cref{thm:bounded-exact-regularization}, the mass of every bounded $\omega$-plurisubharmonic function is a limit of such masses. Hence the infimum and supremum of $\mathrm{e}^{B_0(G)}Z_G$ over smooth $G$ are $v_-(\omega)$ and $v_+(\omega)$. Since $\nu_G$ and $\sigma_G$ are smooth and strictly positive, the last two inequalities prove sharpness on this smaller domain, including the cases $v_-(\omega)=0$ and $v_+(\omega)=+\infty$. The two equivalences follow.
\end{proof}

The difference does not depend on the smooth positive reference volume. For smooth real $a$, replacing $\mathrm{d}V$ by $\mathrm{d}V'=\mathrm{e}^a\mathrm{d}V$ changes both $H_0(\nu)$ and $\operatorname{Ent}(\nu)$ by $-\int_Xa\,\mathrm{d}\nu$, by the transformation formula preceding \cref{prop:hermitian-dual-entropy-bounds}, and finiteness of the entropy is unchanged because $a$ is bounded. Similarly, if $h$ is a smooth real function and $\omega+\ddc h$ is also Hermitian, transporting bounded solutions by $u\mapsto u-h$ preserves $B_0$, and hence this difference. Finally, the difference is constant on $\mathcal P_{\mathrm{ent}}(X)$ exactly when $v_-(\omega)=v_+(\omega)$, that is, when all bounded $\omega$-plurisubharmonic functions have the same mass.
